% Work-MAT | 2. rocnik | hodina 18 | revize 2026-11-09
\documentclass[11pt,a4paper]{article}
\input{../../sablony/pripravy/priprava.sty}
\usepackage{tikz}
\usetikzlibrary{arrows.meta}
\nastavPripravu{2. ročník}{Exponenciální rovnice}{18}{Exponenciální rovnice III: vytýkání}
\begin{document}
\poznamka{Zopakovat převod na společný základ z hodiny 17, teprve pak zavést součty exponenciálních výrazů. Celá hodina je věnována pouze vytýkání; substituce přijde až v následujícím bloku. Vybrané typy jsou inspirovány přiloženým přehledem exponenciálních rovnic, ale gradace a většina zadání jsou nově připravené.}
\blok{Krátké opakování předchozí hodiny}
\begin{prikladReseniBox}[title={Příklad 1}]
Vyřeš $16^{x-1}=8^{x+1}$ převedením obou stran na stejný základ.
\tcblower
\[
\begin{aligned}
2^{4x-4}&=2^{3x+3},\\
4x-4&=3x+3,\\
x&=7.
\end{aligned}
\]
$K=\{7\}$.
\end{prikladReseniBox}
\blok{Když se sčítají mocniny téhož základu}
Dosud jsme upravovali hlavně \emph{součiny} mocnin. V rovnici $2^{x+1}+2^x=24$ ale členy nemůžeme spojit sečtením exponentů. Všimneme si místo toho společného činitele $2^x$ a vytkneme jej, stejně jako při práci s algebraickými výrazy.
\begin{zadaniBox}[title={Věta - Vytýkání stejné exponenciální mocniny}]
Pro libovolný kladný základ $a$ a reálná čísla $p,q$ platí
\[
c\,a^{x+p}+d\,a^{x+q}=a^x\bigl(ca^p+da^q\bigr).
\]
Jde o běžné vytýkání společného činitele; před případným dělením výrazu v závorce ověříme, že není nulový.
\end{zadaniBox}
\begin{prikladReseniBox}[title={Příklad 2}]
Vyřeš $2^{x+1}+2^x=24$.
\tcblower
\[
\begin{aligned}
2\cdot2^x+2^x&=24,\\
(2+1)\cdot2^x&=24,\\
2^x&=8=2^3,\\
x&=3.
\end{aligned}
\]
Zkouška: $16+8=24$; $K=\{3\}$.
\end{prikladReseniBox}
\begin{prikladReseniBox}[title={Příklad 3}]
Vyřeš rovnici ze skupiny úloh na vytýkání:
\[3^x+3^{x+1}=108.\]
\tcblower
\[
\begin{aligned}
3^x+3\cdot3^x&=108,\\
4\cdot3^x&=108,\\
3^x&=27=3^3,\\
x&=3.
\end{aligned}
\]
$K=\{3\}$.
\end{prikladReseniBox}
\begin{prikladReseniBox}[title={Příklad 4}]
Vyřeš $3^{x+2}-3^x=72$.
\tcblower
\[
\begin{aligned}
9\cdot3^x-3^x&=72,\\
8\cdot3^x&=72,\\
3^x&=9=3^2,\\
x&=2.
\end{aligned}
\]
Zkouška: $81-9=72$.
\end{prikladReseniBox}
\begin{prikladReseniBox}[title={Příklad 5}]
Vyřeš rovnici s více různými posuny exponentu:
\[2^{x+1}+2^{x-1}+2^{x+3}=\frac{21}{8}.\]
\tcblower
\[
\begin{aligned}
\left(2+\frac{1}{2}+8\right)2^x&=\frac{21}{8},\\
\frac{21}{2}\cdot2^x&=\frac{21}{8},\\
2^x&=\frac{1}{4}=2^{-2},\\
x&=-2.
\end{aligned}
\]
Všechny tři členy jsme převedli na násobky $2^x$.
\end{prikladReseniBox}
\begin{prikladReseniBox}[title={Příklad 6}]
Vyřeš $2^{x+2}-2^{x-1}=28$ a vysvětli, proč není správné odečítat exponenty v součtu.
\tcblower
\[
\begin{aligned}
\left(4-\frac{1}{2}\right)2^x&=28,\\
\frac{7}{2}\,2^x&=28,\\
2^x&=8,\\
x&=3.
\end{aligned}
\]
Při odčítání dvou mocnin nelze spojovat exponenty jako při dělení; vytkneme společnou mocninu.
\end{prikladReseniBox}
\poznamka{Zdůraznit chybu $a^u+a^v
e a^{u+v}$. Při vytýkání je vhodné psát si nejprve $a^{x+k}=a^k\cdot a^x$, až potom závorku. Nenazývat tento postup substitucí.}
\begin{zadaniReseniBox}[title={Samostatná práce 1}]
\begin{enumerate}[label=\alph*)]
\item $2^{x+3}+2^x=36$.
\item $5^{x+1}-5^x=100$.
\item $4^{x+1}+4^x=80$.
\item $3^{x+2}+3^x=90$.
\item $2^{x+4}-2^{x+1}=112$.
\item $7^{x+1}-6\cdot7^x=49$.
\item $3^{x+1}+3^{x-1}=30$.
\item $2^{x+2}+2^{x-2}=17$.
\item $5^{x+2}-5^{x-1}=124$.
\item $2^{x+1}-2\cdot2^x=0$.
\end{enumerate}
\tcblower
\begin{enumerate}[label=\alph*)]
\item $9\cdot2^x=36$, tedy $x=2$.
\item $4\cdot5^x=100$, tedy $x=2$.
\item $5\cdot4^x=80$, tedy $x=2$.
\item $10\cdot3^x=90$, tedy $x=2$.
\item $14\cdot2^x=112$, tedy $x=3$.
\item $7^x=49$, tedy $x=2$.
\item $\frac{10}{3}\,3^x=30$, tedy $x=2$.
\item $\frac{17}{4}\,2^x=17$, tedy $x=2$.
\item $\frac{124}{5}\,5^x=124$, tedy $x=1$.
\item $0=0$ pro všechna $x$: $K=\mathbb R$.
\end{enumerate}
\end{zadaniReseniBox}
\begin{zadaniReseniBox}[title={Rozšíření 1}]
\begin{enumerate}[label=\alph*)]
\item Rozhodni, zda má rovnice $5^{x+2}-25\cdot5^x=7$ řešení.
\item Co se stane, změní-li se pravá strana na nulu?
\end{enumerate}
\tcblower
\begin{enumerate}[label=\alph*)]
\item Levá strana je $5^x(25-25)=0$. Rovnice $0=7$ nemá řešení.
\item Rovnice $0=0$ platí pro všechna $x\in\mathbb R$.
\end{enumerate}
\end{zadaniReseniBox}
\begin{zadaniReseniBox}[title={Domácí úkol}]
\begin{enumerate}[label=\alph*)]
\item Vyřeš $2^{x+1}+3\cdot2^x=40$.
\item Vyřeš $7^{x+1}-6\cdot7^x=49$.
\item Vyřeš $3^x+3^{x+1}=108$.
\item Vyřeš $2^{x+2}-2^{x-1}=28$.
\item Vyřeš $4^{x+1}+4^x=80$.
\item Vyřeš $3^{x+2}-3^x=72$.
\item Vyřeš $2^{x+3}+2^{x+1}=80$.
\item Vyřeš $5^{x+1}+5^{x-1}=26$.
\item Vyřeš $2^{x+1}+2^{x-1}+2^{x+3}=\frac{21}{8}$.
\item Na konkrétním příkladu vysvětli, proč $2^x+2^x\ne2^{2x}$ obecně.
\end{enumerate}
\tcblower
\begin{enumerate}[label=\alph*)]
\item $5\cdot2^x=40$, tedy $x=3$.
\item $7^x=49$, tedy $x=2$.
\item $4\cdot3^x=108$, tedy $x=3$.
\item $\frac{7}{2}\,2^x=28$, tedy $x=3$.
\item $5\cdot4^x=80$, tedy $x=2$.
\item $8\cdot3^x=72$, tedy $x=2$.
\item $10\cdot2^x=80$, tedy $x=3$.
\item $\left(5+\frac{1}{5}\right)5^x=26$, tedy $5^x=5$ a $x=1$.
\item $\frac{21}{2}\cdot2^x=\frac{21}{8}$, tedy $x=-2$.
\item Například pro $x=2$: $4+4=8$, ale $2^4=16$.
\end{enumerate}
\end{zadaniReseniBox}
\poznamka{Domácí úkol asi 20--30 minut; při nedostatku času lze vybrat prvních sedm a zbylé úlohy ponechat dobrovolně. Všechna řešení používají pouze vytýkání.}
\end{document}
